\section{Continuity and construction of the logarithmic evolution}
\label{sec:evolution}

In this section, we construct the evolution of the normalized log-rate measure.
We use continuity of the averaged generator with respect to the current distribution and an explicit discrete approximation that preserves positivity and total mass.

For \(B>0\) and \(F\in\mathcal P(\R)\), retain the coefficients of Section~\ref{sec:generator} and the relation \(U=B\exp_*F\).
Recall that
\begin{equation}
 \mathcal G_{B,F}\varphi(y)
 =a_{B,F}(y)\varphi'(y)
  +\int_{\R}\bigl[
       \varphi(y+v)-\varphi(y)-v\varphi'(y)
    \bigr]k_{B,F}(y,v)\,\nu_0(\dd v),
 \qquad
 \nu_0(\dd v)=\frac{\dd v}{4\sinh^2(v/2)}.
 \label{eq:evolution-generator}
\end{equation}
The compensation in this expression is over all jump sizes.
Lemma~\ref{lem:bounds} gives, for every compact interval \(I\subset(0,\infty)\), a constant \(C_I<\infty\) such that
\begin{equation}
 0\le k_{B,F}(y,v)\le1,\qquad
 |a_{B,F}(y)-y|\le C_I,\qquad
 m_2:=\int_{\R}v^2\,\nu_0(\dd v)<\infty,
 \qquad B\in I.
 \label{eq:evolution-uniform-bounds}
\end{equation}
These bounds hold for every probability \(F\), without a moment assumption.
We first prove the continuity needed to pass through the nonlinear dependence on \(F\), and then construct a positive weak evolution.

\subsection{Weak continuity of the averaged generator}

\begin{lemma}[Joint continuity of the averaged generator]\label{lem:continuity}
Let \(B_n,B>0\), \(F_n,F\in\mathcal P(\R)\), and \(y_n,y\in\R\).
Suppose \(B_n\to B\), \(F_n\Rightarrow F\), and \(y_n\to y\).
Then
\[
 a_{B_n,F_n}(y_n)\longrightarrow a_{B,F}(y).
 \]
For every \(\varphi\in C_c^2(\R)\),
\begin{equation}
 \mathcal G_{B_n,F_n}\varphi(y_n)
   \longrightarrow\mathcal G_{B,F}\varphi(y).
 \label{eq:generator-joint-continuity}
\end{equation}
Moreover, \(\mathcal G_{B,F}\varphi(y)\) is bounded uniformly over all probabilities \(F\), all \(y\in\R\), and \(B\) in a fixed compact positive interval.
Consequently the scalar functional
\begin{equation}
 H_\varphi(B,F):=F(\mathcal G_{B,F}\varphi)
 \label{eq:averaged-generator-functional}
\end{equation}
is continuous for the usual topology on \(B>0\) and the narrow topology on probabilities \(F\).
It is uniformly bounded when \(B\) is restricted to a compact positive interval.
\end{lemma}

\begin{proof}
\emph{Step 1: coupling the Dirichlet probabilities.}
We use the fixed probability space and the quantile maps of Lemma~\ref{lem:parameterized-dirichlet}.
For a probability \(F\) on \(\R\), put \(Z_j=q_F(S_j)\) and \(V_j=1-T_j^{1/B}\).
Thus the locations have law \(F\), the break fractions have law \(\operatorname{Beta}(1,B)\), and the two sequences are independent.
On the common full-probability event in that lemma, the resulting stick-breaking probability has the form
\begin{equation}
 W_j=V_j\prod_{i<j}(1-V_i),
 \qquad
 Q=\sum_{j\ge1}W_j\delta_{Z_j}.
 \label{eq:evolution-stick-breaking}
\end{equation}
The remaining mass \(R_m=\prod_{j\le m}(1-V_j)\) decreases to zero on the common full-probability event specified in that lemma, and \(Q\) has law \(\DP(BF)\), including for atomic base probabilities.

For the sequence \((B_n,F_n)\), set \(Z_{n,j}=q_{F_n}(S_j)\).
Then \(Z_{n,j}\to Z_j\) almost surely for every \(j\).
For completeness, convergence of the inverse distribution functions holds at every continuity point of the limiting inverse: it follows by bracketing that inverse between continuity points of the limiting distribution function and using \(F_n\Rightarrow F\).
The exceptional set has Lebesgue measure zero because an inverse distribution function is monotone.
Also use common uniforms \(T_j\) to set
\[
 V_{n,j}=1-T_j^{1/B_n},
 \qquad V_j=1-T_j^{1/B}.
 \]
Every fixed finite set of weights and locations then converges almost surely.
On that event, for a bounded continuous function \(f\), the difference between the two stick-breaking integrals is bounded by the difference of their first \(m\) terms plus \(\|f\|_\infty(R_{n,m}+R_m)\).
First letting \(n\to\infty\), and then \(m\to\infty\), proves
\[
 Q_n\Rightarrow Q\quad\hbox{almost surely}.
 \]
The event just used concerns only the countably many weights and locations, so it gives weak convergence of the random measures, not merely convergence for one chosen test function.

Use the independent uniform \(T_*\) from the same probability space and put
\[
 Z_n^*=1-T_*^{1/B_n},\qquad Z^*=1-T_*^{1/B}.
 \]
By the posterior coupling \eqref{eq:palm-coupling}, the log-rate probabilities
\[
 \widehat P_n=(1-Z_n^*)Q_n+Z_n^*\delta_{y_n},
 \qquad
 \widehat P=(1-Z^*)Q+Z^*\delta_y
 \]
have laws \(\DP(B_nF_n+\delta_{y_n})\) and \(\DP(BF+\delta_y)\), respectively.
They converge weakly almost surely.
Their rate-coordinate pushforwards \(P_n=\exp_*\widehat P_n\) and \(P=\exp_*\widehat P\) are the posterior probabilities appearing in the coefficients.

\emph{Step 2: weak-star convergence of the boundary phases.}
For each coupled sample, set \(b_n=e^{y_n}\), \(b=e^y\), and
\begin{equation}
 \begin{split}
 R_n(s)&=b_nM_{P_n}(b_ns)
       =\int_{\R}\frac{\widehat P_n(\dd z)}
                           {s+e^{z-y_n}},\\
 \eta_n(u)&=\xi_{P_n}(b_nu),\qquad
 R(s)=\int_{\R}\frac{\widehat P(\dd z)}
                         {s+e^{z-y}},\qquad
 \eta(u)=\xi_P(bu).
 \end{split}
 \label{eq:scaled-posterior-transform}
\end{equation}
For fixed \(s>0\), the integrand defining \(R_n(s)\) is bounded and continuous in \(z\).
Its derivative with respect to \(y_n\) has absolute value at most \(1/(4s)\), uniformly in \(z\).
Hence \(R_n(s)\to R(s)>0\).
The anchored phase representation \eqref{eq:phase-anchor}, after scaling the rate variable, gives
\begin{equation}
 \log R_n(s)-\log R_n(1)
   =\int_0^\infty\eta_n(u)
       \left(\frac1{s+u}-\frac1{1+u}\right)\dd u,
 \qquad 0\le\eta_n\le1.
 \label{eq:scaled-phase-anchor}
\end{equation}

These identities imply \(\eta_n\to\eta\) weak-star in \(L^\infty(0,\infty)\); explicitly,
\[
 \int_0^\infty f(u)\eta_n(u)\,\dd u
 \longrightarrow\int_0^\infty f(u)\eta(u)\,\dd u
 \qquad\text{for every }f\in L^1(0,\infty).
\]
We prove this assertion by compactness and uniqueness.
Since \(L^1(0,\infty)\) is separable, any subsequence has a further subsequence whose integrals converge on a countable dense set of \(L^1\) tests.
The uniform bound one extends the limits to a bounded linear functional on \(L^1\), represented by a function \(\eta_*\) with \(0\le\eta_*\le1\).
The kernels in \eqref{eq:scaled-phase-anchor} belong to \(L^1\).
Passing to the limit shows that \(\eta_*\) represents \(\log R(s)-\log R(1)\).
Uniqueness of the exponential Stieltjes phase representation, equivalently that for the reciprocal complete Bernstein function, gives \(\eta_*=\eta\) almost everywhere; see \cite[Theorems~6.10 and~7.3]{SSV2010}.
Every subsequential limit is therefore the same, which proves convergence against every \(L^1\) kernel.
No pointwise convergence of the boundary phases has been used.
In particular, exceptional values in their measurable versions do not affect the argument, including after the positive rate scaling.

\emph{Step 3: passing through the posterior averages.}
The drift formula \eqref{eq:log-drift} takes the form
\begin{equation}
 a_{B_n,F_n}(y_n)
 =y_n-\psi_0(B_n+1)-1
   +\E\left[-\log R_n(1)
           +\int_0^\infty\eta_n(u)K(u)\,\dd u\right].
 \label{eq:continuity-drift-formula}
\end{equation}
Here \(\psi_0\) is the digamma function and \(K\in L^1(0,\infty)\) is the correction kernel in \eqref{eq:K}.
The coefficient \(\psi_0(B_n+1)\) retains the original mass parameter \(B_n\) under posterior averaging.
The phase integrals converge almost surely by weak-star convergence and are bounded in absolute value by \(\|K\|_1\).
Moreover,
\[
 \frac{Z_n^*}{2}\le R_n(1)\le1.
 \]
Choose \(B_*<\infty\) with \(B_n\le B_*\) for all \(n\).
The common-uniform coupling gives the integrable bound
\[
 0\le-\log R_n(1)
 \le\log2-\log(1-T_*^{1/B_*}).
 \]
Integrability follows from the logarithmic moment of a \(\operatorname{Beta}(1,B_*)\) variable.
Dominated convergence in \eqref{eq:continuity-drift-formula} proves continuity of the drift.

For the jump part, put \(u=e^v\).
The integral against a single posterior phase has kernel
\begin{equation}
 J_\varphi(y,u)
 =\frac{\varphi(y+\log u)-\varphi(y)
              -\log u\,\varphi'(y)}{(u-1)^2}.
 \label{eq:continuity-jump-kernel}
\end{equation}
Taylor's theorem gives
\[
 |J_\varphi(y,u)|
 \le\frac{\|\varphi''\|_\infty}{2}
                 \frac{(\log u)^2}{(u-1)^2}.
 \]
The dominating function is integrable on \((0,\infty)\): at \(u=1\) the quotient has a removable finite limit, and at zero and infinity it is bounded by integrable multiples of \((\log u)^2\) and \((\log u)^2/u^2\), respectively.
Consequently \(J_\varphi(y_n,\cdot)\to J_\varphi(y,\cdot)\) in \(L^1\).
This convergence and the weak-star convergence of \(\eta_n\) imply convergence of the random jump integrals.
Their uniform bound permits taking expectations.
This proves \eqref{eq:generator-joint-continuity}.

\emph{Step 4: integrating against the varying log-rate law.}
For \(\varphi\in C_c^2(\R)\),
\begin{equation}
 |\mathcal G_{B,F}\varphi(y)|
 \le (|y|+C_I)|\varphi'(y)|
       +\frac{m_2}{2}\|\varphi''\|_\infty.
 \label{eq:compact-test-generator-bound}
\end{equation}
Its right side is bounded uniformly in \(y\).
Joint continuity implies that \(\mathcal G_{B_n,F_n}\varphi\to\mathcal G_{B,F}\varphi\) uniformly on compact \(y\)-sets: otherwise a sequence of points violating uniform convergence has a convergent subsequence contradicting \eqref{eq:generator-joint-continuity}.
Tightness of \(F_n\), this local uniform convergence, the uniform global bound, and weak convergence against the bounded continuous limiting function give \(H_\varphi(B_n,F_n)\to H_\varphi(B,F)\).
\end{proof}

\subsection{A positive Euler construction}

\begin{theorem}[Existence of a log-rate evolution]\label{thm:evolution}
Let \(B_0>0\), \(F_0\in\Ptwo\), and \(T>0\).
Set \(B_t=B_0e^{-t}\).
There is a narrowly continuous curve \((F_t)_{0\le t\le T}\) of probabilities in \(\Ptwo\) such that
\begin{equation}
 \sup_{0\le t\le T}F_t(y^2)<\infty
 \label{eq:evolution-second-moment}
\end{equation}
and
\begin{equation}
 F_t(\varphi)-F_0(\varphi)
   =\int_0^t F_r(\mathcal G_{B_r,F_r}\varphi)\,\dd r,
 \qquad
 0\le t\le T,\quad \varphi\in C_c^2(\R).
 \label{eq:evolution-weak-equation}
\end{equation}
\end{theorem}

Equation~\eqref{eq:evolution-weak-equation} is the meaning of weak solution used here.
The theorem asserts existence; identification of the associated GGC distributions is the subject of the next section.

\begin{proof}
All constants below may depend on \(B_0,T\) and \(F_0(y^2)\), but not on the mesh or on the intermediate probability measures.
The mass parameter lies in the fixed compact interval \(I=[B_0e^{-T},B_0]\).

\emph{Step 1: a positive probability kernel.}
Take \(h=T/N\) and \(\epsilon=\sqrt h\), with \(h\) sufficiently small.
Freeze \(B,F\) during one step, and abbreviate
\begin{equation}
 \begin{split}
 \lambda_\epsilon(y)
   &=\int_{|v|>\epsilon}k_{B,F}(y,v)\,\nu_0(\dd v),\\
 m_\epsilon(y)
   &=\int_{|v|>\epsilon}v k_{B,F}(y,v)\,\nu_0(\dd v),\\
 a_\epsilon(y)&=a_{B,F}(y)-m_\epsilon(y),
 \qquad
 p_\epsilon(y)=1-h\lambda_\epsilon(y).
 \end{split}
 \label{eq:euler-truncated-coefficients}
\end{equation}
The universal dominating measure gives, for \(0<\epsilon\le1\),
\begin{equation}
 \begin{split}
 \lambda_\epsilon(y)
 &\le\int_{|v|>\epsilon}\nu_0(\dd v)
   =\frac{2}{e^\epsilon-1}\le\frac2\epsilon,\\
 |m_\epsilon(y)|
 &\le\int_{|v|>\epsilon}|v|\,\nu_0(\dd v)\\
 &=2\left\{\frac{\epsilon}{e^\epsilon-1}
                      -\log(1-e^{-\epsilon})\right\}
   \le C(1+|\log\epsilon|).
 \end{split}
 \label{eq:euler-truncation-bounds}
\end{equation}
The equalities follow by integrating on the positive half-line and using the evenness of \(\nu_0\); on that half-line its density is \(-\frac{\dd}{\dd v}(e^v-1)^{-1}\).
In particular,
\begin{equation}
 |a_\epsilon(y)|
 \le |y|+C_I+C(1+|\log\epsilon|).
 \label{eq:euler-drift-bound}
\end{equation}
If \(h\le1/16\), then \(p_\epsilon(y)\ge1-2\sqrt h\ge1/2\), uniformly in \(B,F,y\).
Define
\begin{equation}
 \begin{split}
 \Pi_h^{B,F}(y,\dd z)
 ={}&p_\epsilon(y)\,
      \delta_{\,y+h a_\epsilon(y)/p_\epsilon(y)}(\dd z)\\
 &+h\int_{|v|>\epsilon}
       \delta_{y+v}(\dd z)\,k_{B,F}(y,v)\,\nu_0(\dd v).
 \end{split}
 \label{eq:positive-euler-kernel}
\end{equation}
Measurability of the coefficients makes this a Borel kernel.
Both terms are nonnegative and their total mass is exactly one.
Thus $\Pi_h^{B,F}(y,\cdot)$ is a probability for every $y$.
For a test $\varphi$ and a probability $F$, our kernel notation is
\[
 \Pi_h^{B,F}\varphi(y)=\int\varphi(z)\Pi_h^{B,F}(y,\dd z),
 \qquad
 (F\Pi_h^{B,F})(A)=\int\Pi_h^{B,F}(y,A)F(\dd y).
\]
Starting from \(F_0^h=F_0\), we can therefore define recursively
\begin{equation}
 F_{j+1}^h
   =F_j^h\Pi_h^{B_{jh},F_j^h},
 \qquad j=0,\ldots,N-1.
 \label{eq:positive-euler-recursion}
\end{equation}
All coefficients are defined for every probability, so this recursion does not presume existence of the desired solution.

For one step from \(y\), write \(\Delta=z-y\).
Direct calculation from \eqref{eq:positive-euler-kernel} gives
\begin{equation}
 \begin{split}
 \E(\Delta\mid y)&=h a_{B,F}(y),\\
 \E(\Delta^2\mid y)
   &=h\int_{|v|>\epsilon}v^2 k_{B,F}(y,v)\,\nu_0(\dd v)
       +\frac{h^2a_\epsilon(y)^2}{p_\epsilon(y)}.
 \end{split}
 \label{eq:euler-exact-moments}
\end{equation}
Thus the mean compensation is exact.
In particular we do not bound the mean by the divergent untruncated first absolute jump moment.

\emph{Step 2: uniform second moments.}
Put \(M_j^h=F_j^h(y^2)\).
From \eqref{eq:evolution-uniform-bounds},
\[
 2y a_{B,F}(y)\le3y^2+C_I^2.
 \]
Using \(p_\epsilon\ge1/2\), \eqref{eq:euler-drift-bound}, and \eqref{eq:euler-exact-moments}, we obtain
\[
 M_{j+1}^h
 \le (1+3h+C h^2)M_j^h
       +C h+C h^2(1+|\log h|^2).
 \]
Because \(h(1+|\log h|^2)\) is bounded for sufficiently small \(h\), this implies
\begin{equation}
 M_{j+1}^h\le(1+C_T h)M_j^h+C_T h,
 \qquad
 \sup_{\substack{h\ \mathrm{small}\\0\le j\le N}}M_j^h
       \le C_{T,B_0,F_0}<\infty.
 \label{eq:euler-uniform-moments}
\end{equation}
The second inequality follows by iterating the scalar recursion, or by summing its geometric series.
This also proves by induction that each discrete law has a finite second moment, so every moment calculation just used is justified.

\emph{Step 3: compactness in time and in the state variable.}
We realize a finite Markov chain with initial law \(F_0\) and the kernels in \eqref{eq:positive-euler-recursion}.
Only finitely many iterated probability kernels are needed for this realization.
Writing \(Y_j\) for its state at step \(j\), let
\[
 D_{j+1}
  =Y_{j+1}-Y_j-h a_{B_{jh},F_j^h}(Y_j).
 \]
With respect to the filtration $\mathcal F_j=\sigma(Y_0,\ldots,Y_j)$, these are square-integrable martingale differences, by \eqref{eq:euler-exact-moments}.
The same formula and \eqref{eq:euler-uniform-moments} yield
\[
 \E D_{j+1}^2\le C_T h,
 \qquad
 \E\bigl|a_{B_{jh},F_j^h}(Y_j)\bigr|^2\le C_T.
 \]
For \(i<j\), martingale orthogonality gives \(\E|\sum_{\ell=i}^{j-1}D_{\ell+1}|^2\le C_T(j-i)h\).
Cauchy--Schwarz bounds the squared drift sum by \(C_T((j-i)h)^2\).
Consequently
\begin{equation}
 \E|Y_j-Y_i|^2
   \le C_T\bigl\{(j-i)h+((j-i)h)^2\bigr\}.
 \label{eq:euler-time-moment}
\end{equation}

We interpolate the laws rather than the sample paths: for \(t=(j+\theta)h\), \(0\le\theta\le1\), set
\[
 \widetilde F_t^h=(1-\theta)F_j^h+\theta F_{j+1}^h.
 \]
Also let \(\overline F_t^h=F_j^h\) for \(jh\le t<(j+1)h\), with both interpolations equal to \(F_N^h\) at \(T\).
Use the bounded Lipschitz metric
\[
 d_{\mathrm{BL}}(F,G)
 =\sup_{\substack{\|f\|_\infty\le1\\ \operatorname{Lip}(f)\le1}}
                  |F(f)-G(f)|,
 \]
where $\operatorname{Lip}(f)=\sup_{x\ne y}|f(x)-f(y)|/|x-y|$.
This metric metrizes narrow convergence of probabilities on \(\R\).
The coupling in \eqref{eq:euler-time-moment}, followed by Cauchy--Schwarz, gives the uniform estimates
\begin{equation}
 \begin{split}
 d_{\mathrm{BL}}(\widetilde F_t^h,\widetilde F_s^h)
      &\le C_T\sqrt{|t-s|+h},\\
 \sup_{0\le t\le T}
 d_{\mathrm{BL}}(\widetilde F_t^h,\overline F_t^h)
      &\le C_T\sqrt h .
 \end{split}
 \label{eq:euler-time-modulus}
\end{equation}
For the first estimate, compare each interpolated law to a neighboring mesh law and apply \eqref{eq:euler-time-moment}; the squared time increment is absorbed into the constant depending on \(T\).
The second is the adjacent-step case.

The moment bound \eqref{eq:euler-uniform-moments} holds also for the interpolated laws and gives \(\widetilde F_t^h(|y|>R)\le C_{T,B_0,F_0}/R^2\).
They are therefore uniformly tight.
Take a sequence of meshes tending to zero.
At each time in a countable dense subset of \([0,T]\) containing its endpoints, tightness gives a weakly convergent subsequence; diagonal selection gives one subsequence working at all these times.
On the real line this selection can equivalently be obtained from monotonicity of distribution functions and the displayed uniform tail bound.

The first estimate in \eqref{eq:euler-time-modulus} is an asymptotic, uniform time modulus.
Its limit on the dense time set is at most \(C_T\sqrt{|t-s|}\).
Uniform tightness and this bound extend the limiting probabilities uniquely to every time, giving a narrowly continuous curve \(F_t\).
Comparing any time to a finite dense time grid, and then using \eqref{eq:euler-time-modulus}, shows that along the chosen subsequence
\begin{equation}
 \sup_{0\le t\le T}
 d_{\mathrm{BL}}(\widetilde F_t^h,F_t)\longrightarrow0,
 \qquad
 \sup_{0\le t\le T}
 d_{\mathrm{BL}}(\overline F_t^h,F_t)\longrightarrow0.
 \label{eq:euler-uniform-narrow-limit}
\end{equation}
The limits remain probability measures: the uniform tail bound precludes loss of mass.
Approximating a bounded continuous function uniformly on a large compact interval by a Lipschitz function, and controlling the complementary tails, also gives uniform-in-time convergence for every fixed bounded continuous test.
Finally, apply this convergence to \(\min\{y^2,R\}\), and then let \(R\to\infty\).
It follows that \(F_t\in\Ptwo\) for every \(t\) and that \eqref{eq:evolution-second-moment} holds.
Only narrow continuity, not continuity in a second-moment transport metric, is asserted.

\emph{Step 4: consistency with the compensated generator.}
Fix \(\varphi\in C_c^2(\R)\).
Taylor expansion at the no-jump location in \eqref{eq:positive-euler-kernel} gives a remainder bounded by
\[
 \frac{h^2\|\varphi''\|_\infty
                      a_\epsilon(y)^2}{2p_\epsilon(y)}.
 \]
The first-order term is \(h a_\epsilon(y)\varphi'(y)\).
Adding the retained jump terms converts it exactly into the truncated, fully compensated operator.
For the missing small jumps,
\[
 \int_{|v|\le\epsilon}v^2\,\nu_0(\dd v)\le2\epsilon,
 \]
because \(4\sinh^2(v/2)\ge v^2\).
We therefore have the pointwise consistency estimate
\begin{equation}
 \left|
 \frac{\Pi_h^{B,F}\varphi(y)-\varphi(y)}h
       -\mathcal G_{B,F}\varphi(y)
 \right|
 \le C_I\|\varphi''\|_\infty
       \left[\epsilon+
              h\{y^2+1+|\log\epsilon|^2\}\right].
 \label{eq:euler-consistency}
\end{equation}
Here and below the value of the constant \(C_I\) may be enlarged; it remains independent of \(F,y,h\).
Integrating against \(F_j^h\), using \eqref{eq:euler-uniform-moments}, and summing the steps yields an error, uniform in terminal time, at most
\begin{equation}
 C_{T,\varphi}\bigl[\sqrt h+h(1+|\log h|^2)\bigr]
   \longrightarrow0.
 \label{eq:euler-total-error}
\end{equation}
More explicitly, let \(\overline B_r^h=B_{jh}\) on the \(j\)-th mesh interval.
Telescoping the discrete equations and, on the final partial interval, taking the same convex combination as in \(\widetilde F_t^h\), gives
\begin{equation}
 \widetilde F_t^h(\varphi)-F_0(\varphi)
  =\int_0^t
       H_\varphi(\overline B_r^h,\overline F_r^h)\,\dd r
       +e_h(t),
 \qquad
 \sup_{0\le t\le T}|e_h(t)|\longrightarrow0 .
 \label{eq:euler-integrated-consistency}
\end{equation}
Thus the same error bound covers terminal times between mesh points.

\emph{Step 5: the nonlinear limit.}
Lemma~\ref{lem:continuity} says that \(H_\varphi\) is continuous and uniformly bounded on \(I\times\mathcal P(\R)\).
Equations \eqref{eq:euler-uniform-narrow-limit} and \(\sup_r|\overline B_r^h-B_r|\to0\) imply
\begin{equation}
 \sup_{0\le r\le T}
 \left|
 H_\varphi(\overline B_r^h,\overline F_r^h)
       -H_\varphi(B_r,F_r)
 \right|\longrightarrow0.
 \label{eq:euler-nonlinear-limit}
\end{equation}
To justify the uniform assertion without assuming compactness of the entire space of probabilities, suppose it fails and choose times \(r_h\) at which the difference is bounded away from zero.
Along a further subsequence, \(r_h\to r\).
Uniform narrow convergence and continuity of \(F_\cdot\) give \(\overline F_{r_h}^h\Rightarrow F_r\), while \(\overline B_{r_h}^h\to B_r\).
Continuity of \(H_\varphi\) gives a contradiction, also using continuity of \(r\mapsto H_\varphi(B_r,F_r)\).

Passing to the limit in \eqref{eq:euler-integrated-consistency}, we obtain \eqref{eq:evolution-weak-equation} for the fixed test \(\varphi\), at every time.
The subsequence used to construct \(F_\cdot\) was selected only through weak compactness, independently of \(\varphi\).
The same argument therefore applies to every \(\varphi\in C_c^2(\R)\), establishing the stated equation on that entire test class.
\end{proof}

\begin{remark}
The resulting rate measures are Thorin-admissible throughout the finite time interval.
Indeed,
\[
 \int_0^\infty\log(1+1/b)\,
                B_t\exp_*F_t(\dd b)
 =B_tF_t\bigl(\log(1+e^{-y})\bigr)
 \le B_t\bigl(\log2+F_t(|y|)\bigr)<\infty.
 \]
The construction thus produces positive finite Thorin measures with controlled logarithmic tails.
Identifying their associated \(\GGC\) laws with a specified transformation of the initial law is a separate assertion; neither that identification nor nonlinear uniqueness has been used here.
\end{remark}
